\section{Radial caps and their overlap margins}
\label{sec:cap-choices}

The overlapping construction in Section~\ref{sec:tetrahedron} used a
cubic addition to a quadratic radial cap. This appendix describes other
caps that preserve coverage and strict saving. A common criterion proves
all of them first. The subsequent identities give alternative proofs and
coordinate forms, useful when a cap is specified by its square or by a
fractional power of the tilt.

Keep the tetrahedron \eqref{eq:tetrahedron}, the functions $G,H,D_0$,
and the upper and lower parametrizations \eqref{eq:upper-sweep} and
\eqref{eq:lower-sweep}. In this appendix replace their common cap by a
continuous nonnegative function $R_\eps$ on $[-1,1]$. Write
\begin{equation}\label{eq:general-cap-cost}
 Q_R(\eps)=\int_{-1}^1
 \frac{R_\eps(r)^2}{2\sqrt{1+\eps^2D_0(r)}}\,dr.
\end{equation}
Whenever the two sweeps cover $K$, Theorem~\ref{thm:radial} gives
$2N$ triangular-base cylinders, for each $N\ge1$, whose total areas
converge to $2hQ_R(\eps)$ as the mesh tends to zero. The actual swept
projections have the same limiting total area by Theorem~\ref{thm:radial},
although their finite areas can differ from the triangular areas. The tilt
is held fixed throughout the refinement.

\subsection{A common criterion and the exact cost}

\begin{proposition}[Four radial cap choices]\label{prop:cap-choices}
For the tetrahedron and paired sweeps just described, each of the
following caps gives coverage for every sufficiently small $\eps>0$:
\begin{align*}
 R_\eps(r)&=1+\eps^2(H(r)+\beta), &&\beta>0,\tag{L}\\
 R_\eps(r)&=1+\eps^2H(r)+c\eps^3, &&c\ge9/2,\tag{C}\\
 R_\eps(r)&=1+\eps^2H(r)+c\eps^{5/2}, &&c>0,\tag{F}\\
 R_\eps(r)&=\sqrt{1+2\eps^2(H(r)+\beta)},&&\beta>0.\tag{S}
\end{align*}
For (C), coverage in fact holds for every $\eps>0$.
Their normalized costs have the respective expansions
\begin{align}
 Q_R(\eps)&=1+(2\beta-1/30)\eps^2+O(\eps^4), &&\text{(L), (S)},
 \label{eq:fixed-cap-cost}\\
 Q_R(\eps)&=1-\eps^2/30+2c\eps^3+O(\eps^4), &&\text{(C)},
 \label{eq:cubic-cap-cost}\\
 Q_R(\eps)&=1-\eps^2/30+2c\eps^{5/2}+O(\eps^4), &&\text{(F)}.
 \label{eq:fractional-cap-cost}
\end{align}
Thus (L) and (S) give strict saving when $0<\beta<1/60$;
(C) and (F) give strict saving for every stated fixed $c$.
In each case one first fixes a tilt giving coverage and strict saving,
and then chooses a finite sufficiently fine partition.
\end{proposition}
\begin{proof}
The root selection in Proposition~\ref{prop:coverage} is independent
of the cap. Suppose first that
$R_\eps(r)=1+\eps^2(H(r)+b_\eps(r))$, where $b_\eps\ge0$.
If both caps fail, the estimates through \eqref{eq:interface-error}
and the ensuing exact identity apply without change: only nonnegativity
of the cap correction was used there. They give
\begin{equation}\label{eq:cap-failure-bound}
 \frac{xq-yp}{\eps}\le H(p)+H(q)+9\eps.
\end{equation}
The failed caps require the left side to exceed
$H(p)+H(q)+b_\eps(p)+b_\eps(q)$. Consequently
\begin{equation}\label{eq:padding-criterion}
 b_\eps(p)+b_\eps(q)\ge9\eps\quad(p,q\in[-1,1])
\end{equation}
is a sufficient coverage condition. Taking $b_\eps=\beta$, $c\eps$,
or $c\sqrt\eps$ proves (L), (C), and (F), with the stated ranges.
The strict inequality furnished by failure permits equality in
\eqref{eq:padding-criterion}.

For (S), put $A=H+\beta$. Rationalizing the square root gives the
exact expression
\[
 R_\eps=1+\eps^2\left(A-
 \frac{2\eps^2A^2}{(\sqrt{1+2\eps^2A}+1)^2}\right).
\]
Since $A$ is bounded and $\beta>0$, its effective padding is at least
$\beta/2$ for all sufficiently small $\eps$. Condition
\eqref{eq:padding-criterion} again applies. In particular the square-root
cap has not been replaced by its Taylor polynomial.

The cost calculation is uniform in $r$. For (L) and (S), respectively,
\[
 R_\eps^2=1+2\eps^2(H+\beta)+O(\eps^4),
 \qquad R_\eps^2=1+2\eps^2(H+\beta).
\]
For (C) or (F), with $\gamma=3$ or $5/2$,
$R_\eps^2=1+2\eps^2H+2c\eps^\gamma+O(\eps^4)$.
Multiply by
$(1+\eps^2D_0)^{-1/2}=1-\eps^2D_0/2+O(\eps^4)$,
divide by two, and integrate. The integrals already evaluated in
Section~\ref{sec:tetrahedron} give
$\int(H-D_0/4)=-1/30$. This proves all three expansions.
For (L) and (S) with $0<\beta<1/60$, and for (C) and (F) with every
stated $c$, the negative quadratic term dominates the higher-order terms.
Fixing such a tilt and then refining the partition gives the finite saving.
\end{proof}

For reference, the following rational fixed margins all give strict
saving. Several appear in the explicit feet and coordinate realizations
of Appendices~\ref{sec:finite-footprints}--\ref{sec:coordinate-realizations}.
Their exact quadratic cost coefficients are
\begin{center}
\begin{tabular}{c|rrrrrrr}
$\beta$&$1/120$&$1/125$&$1/500$&$4/625$&$1/100$&$1/200$&$1/250$\\
\hline
$2\beta-1/30$&$-1/60$&$-13/750$&$-11/375$&$-77/3750$&$-1/75$&$-7/300$&$-19/750$
\end{tabular}
\end{center}
The square-root choices $\beta=1/120,1/125$ have the same displayed
quadratic coefficients as their linear counterparts, but their caps and
finite covers are different. Cubic choices $c=5,12,20,64$ contribute,
respectively, $10,24,40,128$ at order $\eps^3$. Fractional choices
$c=1,4\sqrt2,8\sqrt2$ contribute $2,8\sqrt2,16\sqrt2$ at order
$\eps^{5/2}$. Each constant here specifies an actual cap rather than an
identification based only on a leading coefficient.

\subsection{Alternative identities for the overlap}

The coverage and saving for all four caps have now been proved.
The following arguments are alternatives to the common criterion, not
additional requirements. The first uses the monotonicity of $G$ in the
squared angular variable; the second keeps positive weights on the radial
parameters; the third avoids expanding a square-root cap.

\paragraph{Fixed margins and monotone polynomials.}
For the roots in \eqref{eq:root-one}--\eqref{eq:root-two}, the segment
equations are
\[
 y=\rho q-\eps xG(q),\qquad
 x=\sigma p+\eps yG(p),\qquad
 \rho+\sigma=2+\eps(xq-yp).
\]
Under joint cap failure for (L), $\rho,\sigma>1$. Put
$E=\rho+\sigma-2$. Since $|\rho-\sigma|\le E$, substitution gives
\[
 E=\eps(\sigma-\rho)pq+\eps^2\{yqG(p)+xpG(q)\}.
\]
All variables except the radial parameters are bounded independently of
$\eps$. Absorbing $\eps E$ for $\eps<1/2$ shows $E=O(\eps^2)$.
Thus $\rho,\sigma=1+O(\eps^2)$ and $x=p+O(\eps)$,
$y=q+O(\eps)$. A further substitution gives uniformly
\[
 E=\eps^2\{q^2G(p)+p^2G(q)\}+O(\eps^3)
 \le\eps^2(H(p)+H(q))+O(\eps^3),
\]
because
$(p^2-q^2)(G(p)-G(q))=(p^2-q^2)^2/2\ge0$.
The failed caps instead give
$E>\eps^2(H(p)+H(q)+2\beta)$, a contradiction for small $\eps$.
Equivalently, one can argue by compactness without estimating the
uniform cubic constant. If failure occurred along $\eps\downarrow0$,
a subsequence would have $p\to P$, $q\to Q$, $x\to P$, $y\to Q$,
and $t/\eps\to PQ$, by the same root equations and cap inequalities.
The exact identity in the proof of Proposition~\ref{prop:coverage} gives
\[
 \lim\frac{xq-yp}{\eps}=(P^2+Q^2)/2+P^2Q^2.
\]
The required lower limit $H(P)+H(Q)+2\beta$ exceeds this by
$(P^2-Q^2)^2/2+2\beta>0$. This establishes the fixed-margin
coverage directly, including all boundary points through the root lemma.

\paragraph{Positive weighted sums and fractional margins.}
The preceding segment equations also give the exact identity
\begin{equation}\label{eq:weighted-radial-identity}
 (1+\eps pq)\rho+(1-\eps pq)\sigma
 =2+\eps^2\{qG(p)y+pG(q)x\}.
\end{equation}
Under failure for (F), the inequalities $\rho,\sigma>1$ imply
$x=p+O(\eps)$ and $y=q+O(\eps)$, uniformly. Thus its right side is
at most $2+\eps^2(H(p)+H(q))+O(\eps^3)$. For $\eps<1$ both
weights on the left are positive. Substituting the two failed caps makes
that side strictly greater than
\[
 2+\eps^2(H(p)+H(q))+2c\eps^{5/2}+O(\eps^3).
\]
The positive fractional term dominates the cubic error. This proves
coverage by a weighted cancellation, without discarding the fractional
term from the cost expansion.

There is a related complementary-coordinate identity. Put
$T=t-\eps xq=1-\rho$ and $S=-t+\eps yp=1-\sigma$. Directly from
the segment equations,
\begin{align}
 (1+\eps xq)T&=t-\eps xy-\eps^2x^2G(q),\label{eq:complementary-one}\\
 (1-\eps yp)S&=-t+\eps xy-\eps^2y^2G(p).\label{eq:complementary-two}
\end{align}
If both caps fail, then $T,S<0$. For $\eps\le1/4$ the two
coefficients are at least $1/2$. Adding the identities therefore gives
\begin{equation}\label{eq:complementary-bound}
 |T|,|S|\le16\eps^2.
\end{equation}
This bounds each missed intercept, rather than only their sum. The
segment equations give $|x-p|,|y-q|\le C\eps$ for a uniform $C$.
For cap (L), substituting the strict failures $T<-\eps^2(H(q)+\beta)$ and
$S<-\eps^2(H(p)+\beta)$ into the sum then requires
\[
 x^2G(q)+y^2G(p)>H(p)+H(q)+2\beta-O(\eps).
\]
The left side is at most $p^2G(q)+q^2G(p)+O(\eps)$ and hence at
most $H(p)+H(q)+O(\eps)$, again a contradiction. Rescaling the axial
coordinate by $t=kz$ and the tilt by $\eps=k\tau$, with
$k=1/\sqrt2$, turns \eqref{eq:complementary-bound} into the explicit
physical-intercept estimate $|T/k|,|S/k|\le16k\tau^2$.

\paragraph{Squared radial parameters.}
The same roots satisfy two further exact identities:
\begin{equation}\label{eq:squared-radial-identities}
 \rho^2=(1-t)^2+2\eps xy+\eps^2x^2,
 \qquad
 \sigma^2=(1+t)^2-2\eps xy+\eps^2y^2.
\end{equation}
For example, expand $(1-t+\eps xq)^2$ and substitute
$2xy=2(1-t)xq+\eps x^2(q^2-1)$ from the first root equation.
For the square-root cap (S), joint failure requires
\[
 2(t/\eps)^2+x^2+y^2>2H(p)+2H(q)+4\beta.
\]
Along any hypothetical sequence of failures approaching zero tilt,
the root and cap inequalities give $x\to P$, $y\to Q$ and
$t/\eps\to PQ$. The limiting right side exceeds the left by
$(P^2-Q^2)^2+4\beta$, a contradiction. For (L), the square of its
cap is at least $1+2\eps^2(H+\beta)$, so the same argument applies.
This retains both the square-root construction and the square-identity
proof for a polynomial cap.


\paragraph{A larger cubic certificate.}
For the cubic choice $c=12$, suppose both selected radial parameters
exceed their caps. The root estimates in Proposition~\ref{prop:coverage}
give
\[
 0<\rho-1,\sigma-1\le\rho+\sigma-2
 =\eps\{(x-p)q-(y-q)p\}\le6\eps^2\le12\eps^2.
\]
Their difference obeys the same bound. Multiplying \eqref{eq:cap-failure-bound} by $\eps^2$ gives an
error $9\eps^3$, which is bounded by the combined padding
$24\eps^3$. Thus this choice also provides a cubic certificate
with that larger constant. Its total limiting cost is
$2\sqrt2(1-\eps^2/30+24\eps^3+O(\eps^4))$.

\subsection{A quantitative complementary calculation at half height}

An exact weighted calculation gives another explicit cubic certificate.
Use coordinates $(X,Y,Z)$ with
$|X|\le1-Z$, $|Y|\le Z$, and their physical realization
$(X/2,Y/2,Z/\sqrt2)$. For $r\in[-1,1]$ define
\[
 \alpha(r)=(1+r^2)/4,\quad \beta(r)=r/2,\quad
 F(r)=r^2(1+r^2)/16.
\]
The two line families are
\begin{align*}
 Y&=ur+\delta X\alpha(r),& Z&=u+\delta X\beta(r),\\
 X&=vs-\delta Y\alpha(s),&1-Z&=v-\delta Y\beta(s).
\end{align*}
The crossing lemma supplies $u,v\ge0$. Both families use the cap
$R_\delta(r)=1/2+\delta^2F(r)+4\delta^3$.
Their exact complementary identity is
\begin{equation}\label{eq:half-height-compensation}
 (u-1/2)(1-\delta Xr)+(v-1/2)(1+\delta Ys)
 =\delta^2\{X^2\alpha(r)+Y^2\alpha(s)\}.
\end{equation}
To verify it, expand
$Z-\delta XY=1/2+(u-1/2)(1-\delta Xr)-\delta^2X^2\alpha(r)$
and its complementary equation, and add.

Suppose both caps fail and $0<\delta\le1/2$. The weights in
\eqref{eq:half-height-compensation} are at least $1/2$, since
$|X|,|Y|\le1$. Its right side is at most $\delta^2$, so each
positive radial excess is at most $2\delta^2$. It follows that
$|Y-r/2|,|X-s/2|\le2\delta^2+\delta/2\le3\delta/2$. Therefore
\[
 X^2\alpha(r)+Y^2\alpha(s)
 \le\frac{s^2\alpha(r)+r^2\alpha(s)}4+\frac94\delta
 \le F(r)+F(s)+\frac94\delta.
\]
Here $|X+s/2|,|Y+r/2|\le3/2$, $\alpha\le1/2$, and the last
inequality is monotonicity in $r^2,s^2$. Failed caps, however, make
\eqref{eq:half-height-compensation} divided by $\delta^2$ exceed
\[
 (1-\delta)(F(r)+F(s)+8\delta)
 \ge F(r)+F(s)+\frac{15}{4}\delta,
\]
because $F\le1/8$. The incompatible constants $9/4$ and $15/4$
prove coverage throughout $0<\delta\le1/2$.

Replacing $4\delta^3$ by $\delta^{5/2}$ gives a second useful
version for all sufficiently small positive $\delta$. To see the fractional cancellation directly, set $a=1/2$,
$P(\xi)=a/2+\xi^2/(2a)$ and $g(\xi)=\xi^2P(\xi)$.
The identity before \eqref{eq:half-height-compensation} becomes
\[
 Z=a+\delta XY-\delta^2P(ar)X^2
       +(u-a)(1-\delta rX).
\]
With both caps failed, $|Y-ar|,|X-as|=O(\delta)$ uniformly and the
positive factors are bounded away from zero. Bounded derivatives of
$P,g$ consequently give
\begin{align*}
 Z&>a+\delta XY+\delta^2(g(Y)-P(Y)X^2)
       +\delta^{5/2}-C\delta^3,\\
 Z&<a+\delta XY-\delta^2(g(X)-P(X)Y^2)
       -\delta^{5/2}+C\delta^3.
\end{align*}
Their incompatibility follows from the exact identity
\[
 g(X)+g(Y)-P(X)Y^2-P(Y)X^2
 =\frac{(X^2-Y^2)^2}{2a}\ge0.
\]
Thus the complementary and the positive-weight arguments both preserve
fractional padding with its proper, non-fourth-order leading remainder.
The affine change $x=2Y$, $y=2X$, $t=2Z-1$ takes these two
half-height families to the earlier canonical families with
$\eps=\delta/2$. The first family becomes the lower family with
radial parameter $2u$, angular parameter $r$, and free coordinate $y=2X$;
the second becomes the upper family with radial parameter $2v$, angular
parameter $s$, and free coordinate $x=2Y$. The cubic and fractional caps become, respectively,
$1+\eps^2H+64\eps^3$ and
$1+\eps^2H+8\sqrt2\eps^{5/2}$.
